\begin{longtable}{p{1.7cm} p{2.3cm} p{1.9cm} p{5.0cm} p{1.1cm} p{1.9cm}}
\caption{Cat\'alogo de requerimientos funcionales (RF)}\label{tab:rf}\\
\toprule
\textbf{ID} & \textbf{T\'itulo} & \textbf{Actor} & \textbf{Descripci\'on} & \textbf{Prioridad} & \textbf{Origen} \\
\midrule
\endfirsthead
\toprule
\textbf{ID} & \textbf{T\'itulo} & \textbf{Actor} & \textbf{Descripci\'on} & \textbf{Prioridad} & \textbf{Origen} \\
\midrule
\endhead
\bottomrule
\endlastfoot
\textbf{RF-01.01} & Recepción y Validación de Cantidades contra OC & Recepcionista de bodega & El sistema debe registrar la recepción física de un proveedor validando las cantidades recibidas por SKU contra la OC y generando una alerta por cada discrepancia, sin bloquear el cierre. & Muy Alta & Cap. 4.1, Anexo A, Cap 12 \\
\textbf{RF-01.02} & Registro Obligatorio de Lote & Recepcionista de bodega & El sistema debe exigir el número de lote del proveedor en campo estructurado para todo SKU con trazabilidad obligatoria, capturándolo por escaneo GS1 (AI 10) o ingreso manual, y bloqueando el avance si el campo está vacío. & Muy Alta & Cap. 4.1, Cap 4.9, Cap. 7.3, \\
\textbf{RF-01.03} & Registro de Fecha de Vencimiento & Recepcionista de bodega & El sistema debe registrar la fecha de vencimiento en formato DD/MM/AAAA para todo SKU con control de vencimiento habilitado, vinculándola al lote y a la ubicación asignada como entrada para la lógica FEFO. & Muy Alta & Cap. 4.1, Cap. 4.8 \\
\textbf{RF-01.04} & Registro de No Conformidades en Recepción & Recepcionista de bodega & El sistema debe permitir registrar una no conformidad durante la recepción seleccionando su tipo desde una lista predefinida e ingresando SKU, lote, cantidad afectada y al menos una fotografía capturada desde el dispositivo. & Alta & Cap. 4.1, Cap. 4.9, Cap 8 (Entrevista de Calidad), RT-17.06 \\
\textbf{RF-01.05} & Cuarentena Automática por No Conformidad & Recepcionista de bodega & Al confirmar una no conformidad, el sistema debe cambiar automáticamente a estado ‘Cuarentena’ la cantidad afectada del SKU/lote correspondiente, bloqueándola para picking, y notificar al perfil de calidad mediante una alerta en la aplicación. & Alta & Cap. 4.1, Cap. 4.9, RT-17.06 \\
\textbf{RF-01.06} & Lectura y Decodificación de Códigos GS1 & Recepcionista de bodega & El sistema debe decodificar códigos de barras GS1 interpretando los AIs 01 (GTIN), 10 (lote), 17 (vencimiento) y 00 (SSCC), autocompletando los campos en pantalla; ante código inválido o no registrado, habilita el ingreso manual. & Alta & RT-17.06, Cap. 16.2 (GS1) \\
\textbf{RF-01.07} & Generación e Impresión de Etiqueta SSCC & Recepcionista de bodega & Al finalizar el registro de un pallet, el sistema debe generar automáticamente un SSCC bajo estándar GS1-128 vinculado al GTIN, lote, fecha de vencimiento y fecha de recepción, e imprimir la etiqueta en la impresora térmica del andén. & Media & Cap. 4.2 , Cap 12 (GS1), RT-05.23 \\
\textbf{RF-01.08} & Asignación de SSCC a Ubicación de Almacenamiento & Operador de bodega & El sistema debe permitir asociar un SSCC a una dirección de almacenamiento (pasillo-columna-nivel) escaneando ambas etiquetas y validando la compatibilidad térmica antes de confirmar. & Media & Cap. 4.2, Cap. 12 (GS1), RT-05.23 \\
\textbf{RF-01.09} & Integración de Recepción con ERP & Recepcionista de bodega & Al cerrar una recepción, el sistema debe enviar automáticamente la transacción de ingreso de inventario al ERP mediante su interfaz de integración; si el ERP no está disponible, encola la transacción para reintento automático. & Muy Alta & Cap. 5, Anexo A, RT-03.13 \\
\textbf{RF-01.10} & Gestión de Fichas de Producto & Administrador de catálogo / Jefe de abastecimiento & El sistema debe permitir crear y mantener fichas de producto con atributos obligatorios: descripción, formato, GTIN, SKU interno, tipo de almacenamiento, temperatura mínima y máxima (\ensuremath{^\circ}C), peso neto (kg), dimensiones (cm), días de vida útil, indicadores de trazabilidad y control de vencimiento (sí/no), clasificación sanitaria y RUT del proveedor. & Muy Alta & Cap 14.1 (pasa de 8.400 a 9.500 SKu), Cap 8 (Hugo), Cap 16.1 (N11) \\
\textbf{RF-01.11} & Validación de Compatibilidad Térmica en Putaway & Operador de bodega & El sistema debe validar la compatibilidad térmica entre el tipo de almacenamiento del SKU y la zona de temperatura de la ubicación destino al realizar un putaway, bloqueando la operación si no son compatibles. & Muy Alta & Cap. 2.3, Cap. 4.9, Cap. 10 (restricción 2) \\
\textbf{RF-02.01} & Configuración de topología de instalaciones & Administrador del sistema / Jefe de bodega & El sistema debe permitir registrar y configurar cada instalación de la red logística (centros de distribución y plataformas de cross-docking), definiendo su topología interna: zonas operativas (seco con rango 15–25 \ensuremath{^\circ}C, refrigerado con rango 2–8 \ensuremath{^\circ}C, congelado con rango $-$18 a $-$25 \ensuremath{^\circ}C), pasillos, racks, niveles y posiciones individuales de almacenamiento. La configuración debe estar disponible en modo offline. & Muy Alta & Cap. 2.3, RT-02.12 \\
\textbf{RF-02.02} & Consolidación de stock físico multi-sitio & Jefe de bodega / Gerente de operaciones & El sistema debe consolidar el stock físico en unidades de cada SKU en todos los nodos de la red logística. & Alta & Cap. 2.3 \\
\textbf{RF-02.03} & Reasignación de ubicaciones (slotting) & Jefe de bodega & El sistema debe calcular periódicamente una propuesta de reasignación de ubicaciones (slotting) basada en la clasificación de rotación (ABC) del último período configurable, considerando además peso, volumen y compatibilidad de zona. Los productos de clasificación A deben proponerse en ubicaciones cercanas a los andenes de salida. La propuesta debe ser aprobada por el jefe de bodega antes de generar instrucciones de movimiento. & Muy Alta & Cap. 8 (Ximena) \\
\textbf{RF-02.04} & Conteo cíclico en modo ciego & Operario de bodega / Jefe de bodega & El sistema debe generar un plan de conteo cíclico mensual que cubra una cantidad de posiciones definidas, seleccionando posiciones por criterios configurables (rotación, valor, discrepancia histórica). Las tareas de conteo deben asignarse a los dispositivos de los operarios de bodega en modo de conteo ciego, sin mostrar la cantidad teórica. Tras el ingreso de la cantidad física, el sistema compara contra el teórico y registra la diferencia con su causa en campo estructurado. & Alta & Cap. 4.2, Cap. 14.1 \\
\textbf{RF-02.05} & Cálculo de stock disponible para venta & Preventista / Jefe de bodega & El sistema debe calcular el stock disponible para venta por SKU y por centro de distribución, definido como la diferencia entre el stock físico registrado y el inventario comprometido por pedidos confirmados pendientes de preparación. En modo desconectado, la consulta debe reflejar el stock al momento de la última sincronización, indicando la antigüedad del dato. & Alta & Cap. 16.1 (Decisión \#8) \\
\textbf{RF-02.06} & Gestión de inventario en tránsito & Jefe de bodega / Operaciones & El sistema debe registrar como «en tránsito» todo inventario despachado desde un nodo de origen hacia un nodo destino, desde la carga confirmada hasta la recepción confirmada en destino, y generar la transacción de transferencia correspondiente en el ERP. & Media & Cap. 2.3 (abastecimiento desde Talca) \\
\textbf{RF-02.07} & Consulta de nivel de ocupación por zona & Jefe de bodega / Gerente de operaciones & El sistema debe permitir consultar el nivel de ocupación de cada centro de distribución, desglosado por zona operativa (seco, refrigerado, congelado), expresado como porcentaje de ubicaciones físicas utilizadas respecto de la capacidad total configurada. &  &  \\
\textbf{RF-02.08} & Gestión FEFO y alertas de vida útil & Jefe de bodega & El sistema debe dirigir la preparación de pedidos (picking) asignando tareas de recolección al dispositivo del preparador, indicando: ubicación, producto, lote, cantidad, y secuencia optimizada de recorrido. El preparador confirma cada línea en el dispositivo. Al no encontrar un producto o encontrar cantidad insuficiente, el preparador debe registrar el faltante con su causa en un campo estructurado (sin stock, producto dañado, producto vencido, error de ubicación, otro). &  &  \\
\textbf{RF-02.09} &  & Jefe de bodega & El sistema debe gestionar el almacenamiento de productos con fecha de vencimiento aplicando la regla FEFO (First Expired, First Out), priorizando en la preparación de pedidos los lotes con fecha de vencimiento más próxima. El sistema debe alertar cuando un producto en stock alcance un umbral de vida útil remanente configurable. &  &  \\
\textbf{RF-02.10} & Ficha de trazabilidad de lote & Perfil de calidad / Jefe de bodega & El sistema debe proveer una interfaz de trazabilidad que permita buscar y detallar el ciclo de vida completo de un lote. Al consultar un lote específico, debe desglosar su estado actual, ubicación física en bodega, datos de recepción de proveedor y detalle de salida (clientes/mermas). & Muy Alta & Cap. 1, Cap. 4.9, Cap. 9.5 \\
\textbf{RF-03.01} & Consulta de stock disponible en línea & Preventista & La app debe mostrar, con conectividad, el stock disponible en tiempo real para cada SKU al tomar el pedido. & Muy Alta &  \\
\textbf{RF-03.02} & Consulta de stock disponible sin conectividad & Preventista & La app debe mostrar, sin conectividad, el stock según la última sincronización válida, con su antigüedad visible. & Muy Alta &  \\
\textbf{RF-03.03} & Reserva de stock al confirmar el pedido & Preventista & El sistema debe reservar el stock comprometido en el momento de la confirmación del pedido en el sistema de gestión, por orden de llegada. & Alta &  \\
\textbf{RF-03.04} & Consulta de crédito y deuda vencida en línea & Preventista & La app debe mostrar, con conectividad, el crédito disponible y la deuda vencida del cliente al tomar el pedido. & Muy Alta &  \\
\textbf{RF-03.05} & Consulta de crédito y deuda vencida sin conectividad & Preventista & La app debe mostrar, sin conectividad, el crédito y la deuda según la última sincronización válida, con su antigüedad visible. & Muy Alta &  \\
\textbf{RF-03.06} & Consulta de historial de compra del cliente & Preventista & La app debe mostrar el historial de compra del cliente (productos y frecuencia). & Alta &  \\
\textbf{RF-03.07} & Aplicación de promociones a una línea de pedido & Preventista & La app debe permitir aplicar una promoción vigente (3x2, descuento por volumen) a una línea de pedido, calculando el descuento automáticamente. & Alta &  \\
\textbf{RF-03.08} & Alerta por superación del límite de crédito & Preventista & El sistema debe alertar al preventista cuando el pedido, sumado a la deuda vencida, supere el límite de crédito del cliente. & Muy Alta &  \\
\textbf{RF-03.09} & Bloqueo de envío por exceso de crédito & Preventista & El sistema debe bloquear el envío del pedido cuando este, sumado a la deuda vencida, supere el límite de crédito en más de [110] \%. & Muy Alta &  \\
\textbf{RF-03.10} & Autorización de excepción por supervisor de crédito & Supervisor de credito & El sistema debe permitir que un supervisor autorice, de forma registrada, el envío de un pedido bloqueado. & Alta &  \\
\textbf{RF-03.11} & Aplicación del precio vigente al despacho por defecto & Administrador comercial & El sistema debe aplicar, por defecto, el precio vigente al despacho cuando el precio cambió desde la toma del pedido. & Alta &  \\
\textbf{RF-03.12} & Configuración de excepción de precio por cliente o cadena & Administrador comercial & El sistema debe permitir configurar, por cliente o cadena, que se respete el precio de toma de pedido en lugar del precio de despacho. &  &  \\
\textbf{RF-03.13} & Alerta al cliente por cambio de precio & Cliente & El sistema debe generar una alerta al cliente cuando el precio de despacho difiera del precio pactado, si no tiene excepción configurada (RF-03.12). &  &  \\
\textbf{RF-03.14} & Información de ventana de entrega comprometible &  &  &  &  \\
\textbf{considerando stock} & capacidad de despacho y restricciones del cliente.,El módulo de planificación de rutas expone capacidad disponible (dependencia externa).,El preventista visualiza la ventana disponible antes de confirmar; la promesa queda registrada.,Muy Alta, &  &  &  &  \\
\textbf{RF-03.15} & Sugerencia de productos al preventista & Preventista & El sistema debe generar sugerencias de productos usando historial de compra, frecuencia, productos habituales y promociones vigentes. & Alta &  \\
\textbf{RF-03.16} & Generación de identificador único de pedido offline & Aplicacion de preventa & La app debe generar un identificador único en el dispositivo al crear un pedido sin conectividad. & Muy Alta &  \\
\textbf{RF-03.17} & Detección y descarte de pedidos duplicados al sincronizar & Preventista & El sistema debe detectar y descartar automáticamente los reenvíos duplicados de un pedido con el mismo identificador único al sincronizar. & Muy Alta &  \\
\textbf{RF-04.01} & Asignación y secuenciación automática de rutas & Planificador & El sistema debe asignar pedidos a la flota, secuenciando la ruta en función de la capacidad volumétrica y zonas logísticas de entrega. & Alta & Cap 4.4, Criterio de aceptación \#7 \\
\textbf{RF-04.02} & Modificación manual de rutas y reasignación de pedidos & Planificador & El sistema debe permitir modificar manualmente la secuencia de ruta y reasignar pedidos entre vehículos & Media &  \\
\textbf{RF-04.03} & Bloqueo de asignación por exceso de capacidad & Planificador / Sistema & El sistema debe bloquear la adición de líneas de pedido a una ruta si la sumatoria excede los límites de peso o volumen nominales del vehículo. & Alta &  \\
\textbf{RF-04.04} & Ajuste de secuencia de ruta por ventanas horarias & Sistema & El sistema debe permitir la configuración de ventanas horarias específicas por cliente y ajustar la secuencia para cumplirlas. & Alta &  \\
\textbf{RF-04.05} & Restricción de vehículos por requisito de cadena de frío & Planificador / Sistema & El sistema debe restringir la asignación de pedidos refrigerados o congelados exclusivamente a vehículos con equipo de frío. & Alta &  \\
\textbf{RF-04.06} & Notificación de hora estimada de llegada al cliente & Sistema & El sistema debe crear un aviso al contacto del cliente con la hora estimada de llegada incluyendo una ventana de 30 minutos. & Media &  \\
\textbf{RF-04.07} & Parametrización de ventanas horarias & Planificador / Sistema & El sistema debe permitir parametrizar franjas horarias de recepción en la ficha maestra de cada cliente (Ventana Horaria). & Alta &  \\
\textbf{RF-04.08} & Cálculo integral del costo de entrega realizada & Sistema & El sistema debe calcular el costo de cada entrega tras haber sido realizada. & Alta &  \\
\textbf{RF-05.01} & Generación de misiones de picking & Jefe de Bodega & El jefe de bodega debe poder iniciar la generación manualmente o configurar la generación automática al cierre de pedidos del turno. Las misiones deben permitir agrupación por pedido individual o por consolidación de pedidos de una misma ruta. & Muy Alta & Cap. 4.5 \\
\textbf{RF-05.02} & Validación de líneas por escaneo GS1 & Preparador & El sistema debe exigir la validación de cada línea recolectada escaneando el código GS1. El escaneo debe validar el SKU y capturar el Lote exacto tomado. Permite ingreso manual solo como contingencia auditable. & Muy Alta & Cap. 4.5, Cap. 12 \\
\textbf{RF-05.03} & Motivo obligatorio de faltante & Preparador & El sistema debe exigir la selección obligatoria de un motivo parametrizable de una lista estructurada, cada vez que se reporte una cantidad recolectada inferior a la solicitada. & Muy Alta & Cap. 8, Crit. \#15 \\
\textbf{RF-05.04} & Secuenciación térmica de misiones & Sistema & El sistema debe secuenciar las líneas de la misión de preparación priorizando las zonas en el siguiente orden: seco $\rightarrow$ refrigerado $\rightarrow$ congelado, para minimizar el tiempo de exposición de los productos de frío fuera de su zona térmica. & Alta & Cap. 4.5 \\
\textbf{RF-05.05} & Asignación de inventario FEFO en picking & Preparador & El sistema debe asignar el inventario para la preparación de pedidos aplicando estrictamente la regla lógica FEFO (First Expired, First Out) según la fecha de vencimiento del lote. & Muy Alta & Cap. 4.8, 9.10 \\
\textbf{RF-05.06} & Envío de preparación al ERP &  & El sistema debe enviar al ERP, al cerrar una misión de preparación, el detalle consolidado por cada línea: SKU, cantidad efectivamente recolectada y lote correspondiente, vinculado al ID del pedido original. & Muy Alta & Cap. 4.5 (segundo punto de pérdida) \\
\textbf{RF-05.07} & Carga dirigida inversa a ruta & Cargador / Sistema & El sistema debe dirigir la carga de pedidos preparados en el camión siguiendo la secuencia inversa de la ruta asignada, garantizando que el primer cliente a visitar sea el último cargado (más accesible). El cargador debe confirmar cada bulto/pallet escaneando su SSCC. & Alta & Cap. 4.5, Criterio de Aceptación \#14 \\
\textbf{RF-06.01} & Confirmación de recepción de bultos & Conductor & El sistema debe permitir marcar la recepción conforme de bultos mediante validación por línea de pedido individual o marcado masivo por guía. & Alta &  \\
\textbf{RF-06.02} & Captura de firma digital & Conductor / Cliente & El sistema debe permitir capturar la firma digital del receptor como evidencia de entrega, asociada a la guía de despacho electrónica. & Muy Alta &  \\
\textbf{RF-06.03} & Reporte de rechazo de productos & Conductor & El sistema debe permitir reportar un rechazo total o parcial seleccionando la cantidad y el motivo desde un catálogo predefinido. & Media &  \\
\textbf{RF-06.04} & Reporte de local cerrado y reagendamiento automático & Conductor & El sistema debe permitir reportar un local cerrado, indicando la dirección del cliente. & Media &  \\
\textbf{RF-06.05} & Registro de recaudación de pagos en efectivo & Conductor & El sistema debe permitir digitar el monto exacto recaudado en pagos realizados en efectivo & Muy Alta &  \\
\textbf{RF-06.07} & Sincronización automática de registros de turno & Conductor repartidor, Jefe de bodega / Gerente de operaciones & El sistema debe sincronizarse automaticamente, al detectar la reconexión de señal o el regreso al centro de distribución, sincronizar automaticamente todos los registros generados durante el turno, recuperando la totalidad de los despachos registrados, consolidandolos en un reporte de cierre. & Muy Alta & Cap. 4.6 y 4.7 \\
\textbf{RF-06.08} & Autenticacion conductor externo mediante OTP & Conductor repartidor y Jefe de flota / Transportista & El sistema debe permitir al conductor de una empresa externa iniciar su jornada usando las credenciales temporales por OPT definido en el RF-15.02. sin necesidad de correo corporativo. & Alta & Restricción \#6, RT-12.11, Cap. 16.1 (Decisión \#5) \\
\textbf{RF-06.09} & Captura de devolución de envaces retornables & Conductor repartidor & Al registrar una entrega, el sistema debe permitir al conductor capturar la cantidad de canastillos y pallets devueltos por el cliente, identificando el tipo y estado. & Alta & Cap. 4.8, Cap. 9.7, Criterio Aceptación \#12 \\
\textbf{RF-06.10} & Actualización de saldo de envases retornables & Conductor repartidor & El sistema debe actualizar el saldo del cliente y del parque retornable en linea o en modo diferido según conectividad, asociando el registro a la entrega o a la guia correspondiente & Alta & Cap. 4.8, Cap. 9.7, Criterio Aceptación \#12 \\
\textbf{RF-06.11} & Confirmación de recepción mediante QR & Conductor repartidor & El sistema debe ofrecer al cliente la opción de confirmar la recepción mediante un código QR único generado por la app del conductor, cuando el cliente disponga de telefono inteligente. & Media & Cap. 16.1 (Decisión \#15) \\
\textbf{RF-06.12} & Confirmacion de recepción sin conectividad & Cliente , Conductor repartidor & El sistema debe permitir confirmar la recepción del cliente aun si este no dispone de telefono o conectividad, mediante captura de firma en pantalla o el registro fotografico de la recepcion, dejando evidencia vinculada a la guia de despacho electrónica en su acuse de recibo. & Media & Cap. 16.1 (Decisión \#15) \\
\textbf{RF-06.13} & Impresión de comprobante térmico en terreno. & Conductor repartidor propio y Cliente canal tradicional & La app de reparto debe permitir imprimir en terreno un comprobante térmico de entrega para el cliente usando una impresora portátil Bluetooth en aquellos casos donde el cliente requiera respaldo físico. & Media & RT-17.06 (impresora térmica portátil) \\
\textbf{RF-07.02} & Rendición Digital Individual de Efectivo & Cajero Liquidador / Conductor repartidor propio, Conductor repartidor externo & El sistema debe generar la planilla digital de rendición individual al retorno del camión, contrastando los valores declarados por el conductor en la aplicación móvil contra las entregas al contado realizadas y facturas cobradas. &  & Cap. 4.7, CA-10 \\
\textbf{RF-07.03} & Registro Obligatorio de Causales de Descuadre en Caja & Cajero Liquidador/ Conductor repartidor propio, Conductor repartidor externo & El sistema debe obligar al registro estructurado y categorizado de causales para cualquier diferencia detectada en la rendición de caja (descuentos comerciales autorizados, diferencias de cambio, siniestro en ruta), bloqueando el cierre de la liquidación si existe saldo sin justificar. &  & Cap. 4.7, 9.7, CA-10 \\
\textbf{RF-07.04} & Consola de Gestión de Cartera de Crédito & Ejecutivo de Cobranza/ Gerente de Finanzas, Clientes del canal tradicional, Clientes del canal moderno & El sistema debe proveer una consola de gestión para la cartera de crédito a 30 días sincronizada con el ERP, asignando carteras de clientes a ejecutivos de cobranza y registrando bitácoras de gestión, compromisos de pago y fechas de recontacto. &  & Cap. 4.7 \\
\textbf{RF-07.05} & Alerta y Bloqueo de Crédito en Preventa & Preventista / Ejecutivo de Cobranza, Gerente Comercial & El sistema debe notificar en la aplicación móvil de preventa el estado de bloqueo o mora de un cliente (>30 días o cupo excedido), bloqueando preventivamente el ingreso de pedidos a crédito y permitiendo únicamente venta al contado o cobro de deuda. &  & Cap. 4.3, 8, RT-09.01 \\
\textbf{RF-07.06} & Cobro con Terminal POS Móvil en Ruta & Conductor repartidor propio, Conductor repartidor externo / Clientes del canal tradicional, Clientes del canal moderno, Cajero Liquidador & El sistema debe integrarse vía Bluetooth con terminales de pago portátil (POS) en la aplicación móvil de reparto para cobro con tarjeta en ruta, capturando el código de autorización bancario y conciliando el voucher automáticamente en la liquidación diaria. &  & RT-17.06, Cap. 9.7 \\
\textbf{RF-07.07} & Consulta de Estado de Cuenta en Portal de Clientes & Clientes del canal tradicional, Clientes del canal moderno/ Ejecutivo de Cobranza, Gerente de Finanzas & El sistema debe disponer en el portal web de autoatención el estado de cuenta y la cartola de antigüedad de deuda (tramos corriente, 30, 60, 90+ días) descargable en formatos PDF y Excel para clientes autenticados. &  & Cap. 9.7, RT-16.30 \\
\textbf{RF-07.08} & Validación en Tesorería de Pagos Web & Cajero Liquidador / Clientes del canal moderno, Gerente de Finanzas & El sistema debe proveer una bandeja de validación para pagos reportados en el portal web de clientes, permitiendo al cajero verificar el comprobante bancario y autorizar su traspaso estructurado al ERP. &  & Cap. 4.7, 9.7, RT-16.30 \\
\textbf{RF-07.09} & Interfaz Asíncrona de Cobranzas hacia ERP & Jefe de TI / Gerente de Finanzas, Cajero Liquidador & El sistema debe transferir de forma estructurada, asíncrona e idempotente los comprobantes de recaudación y cierres comerciales aprobados hacia el ERP para actualizar automáticamente las cuentas por cobrar. &  & Cap. 5, Restricción R4 \\
\textbf{RF-07.10} & Registro de Abonos y Pagos Parciales en Entrega & Conductor repartidor propio, Conductor repartidor externo / Clientes del canal tradicional, Cajero Liquidador, Gerente de Finanzas & El sistema debe permitir el registro de anticipos o pagos parciales en entregas al contado desde la aplicación móvil de reparto, emitiendo un comprobante digital de abono y traspasando el saldo insoluto a la cuenta corriente del cliente en el ERP. &  & Cap. 4.7 \\
\textbf{RF-07.11} & Cobranza Presencial de Facturas por Preventistas & Preventista / Cajero Liquidador, Clientes del canal tradicional & El sistema debe permitir al preventista registrar la cobranza en efectivo de facturas vencidas en el local del cliente desde su terminal móvil, emitiendo un recibo provisional digital y forzando su rendición diaria en caja. &  & Cap. 4.3, Cap. 8 \\
\textbf{RF-08.01} & Registro de Devoluciones en Ruta & Conductor repartidor & El sistema debe permitir al conductor registrar la devolución de productos en el momento de la entrega, capturando SKU, cantidad y motivo (daño / vencimiento / producto no pedido / error), operando en modo offline. & Muy Alta & Cap. 4.6, Cap. 4.8, RT-17.01 \\
\textbf{RF-08.02} & Gestión del Destino de Productos Devueltos & Jefe de bodega & Al recepcionar una devolución en el andén del CD, el sistema debe permitir seleccionar el destino final del producto (reingreso a stock / venta con descuento / devolución al proveedor / destrucción) y enviar la transacción al ERP mediante su interfaz de integración. & Muy Alta & Cap. 4.8, Cap 10 (restriccion 4) \\
\textbf{RF-08.03} & Gestión de Cuenta Corriente de Envases Retornables & Conductor repartidor & El sistema debe mantener una cuenta corriente de envases retornables por cliente, actualizando el saldo en cada entrega y devolución registrada por el conductor. & Alta & Cap. 4.8, Cap. 16.1 (Decisión \#10), Cap 9.7 \\
\textbf{RF-08.04} & Alertas por Exceso de Envases Retornables & Preventista & El sistema debe generar una alerta al preventista cuando el saldo de envases retornables de un cliente supere el umbral configurado por el administrador. & Alta & Cap. 4.8, Cap 9.7 \\
\textbf{RF-08.05} & Registro de Mermas por Vencimiento & Preparador de bodega / Preventista & El sistema debe permitir registrar mermas por vencimiento detectadas en conteo cíclico, preparación de pedidos o góndola del cliente, asociando cada registro al SKU y al lote correspondiente. & Alta & Cap. 4.8, Cap 10 (restriccion 4) \\
\textbf{RF-08.06} & Reporte de Pérdida de Envases Retornables & Jefe de operaciones / Finanzas & El sistema debe reportar la pérdida de envases retornables desglosada por cliente, tipo de envase y zona, con actualización diaria. & Alta & Cap. 4.8, Cap 18 (Criterio Aceptación \#12) \\
\textbf{RF-08.07} & Integración de Devoluciones y Mermas al Costo de Servir & Finanzas & El sistema debe incluir el costo de devoluciones y mermas por cliente como componente del cálculo del costo de servir (Épica 11). & Alta & Cap. 9.8 \\
\textbf{RF-09.01} & Trazabilidad de producto foward tracing & Equipo de calidad & El sistema debe permitir la trazabilidad 'hacia adelante' (forward tracing): dado un lote específico, debe listar todos los clientes que recibieron productos de ese lote, con fecha de entrega, guía y cantidad. & Muy Alta & Cap. 1, Cap. 9.6, Criterio Aceptación \#1 \\
\textbf{RF-09.02} & Trazabilidad de producto backward tracing & Equipo de calidad & El sistema debe permitir la trazabilidad 'hacia atrás' (backward tracing): dado un producto en el local del cliente (SKU + lote), debe identificar el origen (proveedor, fecha de recepción y orden de compra) & Muy Alta & Cap. 4.9, Cap. 9.6 \\
\textbf{RF-09.03} & Registro de temperaturas de cámaras de refigerado y congelado & Sensor de temperatura & El sistema debe registrar de forma automática y continua la temperatura de las cámaras de refrigerado y congelado, con lecturas cada 5 minutos como mínimo & Muy Alta & Cap. 4.9, Cap. 7.3, Criterio Aceptación \#3 \\
\textbf{RF-09.04} & Registro de temperaturas de camiones con equipo de frío & Sensor de temperatura & El sistema debe registrar de forma automática y continua la temperatura en los camiones con equipo de frío (propios + de transportistas que cuenten con sensor), durante todo el trayecto desde la carga hasta la entrega. & Muy Alta & Cap. 4.6, Cap. 4.9, Criterio Aceptación \#3 \\
\textbf{RF-09.05} & Alerta de excusiones térmicas & Sistema de monitoreo & El sistema debe generar una alerta en tiempo real (notificación en cabina y en plataforma de monitoreo) cuando se detecte una excursión de temperatura que supere el umbral definido para el tipo de producto. & Muy Alta & Cap. 8 (Katherine), Cap. 16.1 (Decisión \#4) \\
\textbf{RF-09.06} & Parametrización de excusiones térmicas & Jefe de calidad/operaciones & El sistema debe permitir configurar los parámetros que determine una excursión térmica (definida por tiempo/grados) & Muy Alta & Cap. 8 (Katherine vs Nelson), Cap. 16.1 (Decisión \#4) \\
\textbf{RF-09.07} & Bloqueo de despacho por excusiones térmicas & Conductor & El sistema permitira al conductor bloquear el despacho del cambion según su criterio, basandose en el tipo de producto y criterio del mandante & Media & Cap. 8 (Katherine vs Nelson), Cap. 16.1 (Decisión \#4) \\
\textbf{RF-09.08} & Registro de pedido rechazado por cliente & Conductor & El sistema debe permitir al conductor registrar el rechazo de un producto entregado por 'sospecha de ruptura de cadena de frío', asociando el evento al sensor de temperatura del camión en ese tramo. & Muy Alta & Cap. 4.9, Cap. 8 \\
\textbf{RF-09.09} & Registro de lotes con control sanitario & Preparador & El sistema debe registrar, en el momento de la preparación (picking), el lote específico que se le asigna a cada línea de pedido, para los productos con control sanitario. & Muy Alta & Cap. 4.5, Cap. 4.9 \\
\textbf{RF-09.10} & Generación de cumplimiento de cadena de frío &  & El sistema debe generar la evidencia de cumplimiento de cadena de frío exigible ante la autoridad sanitaria y ante el cliente, para un envío determinado. & Muy Alta & Cap. 4.9 ; Cap. 18 criterio N\ensuremath{^\circ}3 \\
\textbf{RF-11.01} & Distribución Automática de Reportes Ejecutivos & Gerente de Finanzas / Gerente Comercial, Jefa de Calidad & El sistema debe programar y distribuir vía correo electrónico reportes ejecutivos consolidados (en formatos PDF y XLSX) con periodicidad semanal y mensual, desglosando la evolución de OTIF, Fill Rate, Merma por vencimiento y Costo de Servir por zona y canal de comercialización. &  & Cap. 7.1, 7.2, 9.8, CA-11 \\
\textbf{RF-11.02} & Cálculo del Costo de Servir por Entrega y Cliente & Gerente de Finanzas / Gerente Comercial & El sistema debe calcular el Costo de Servir real por entrega y por cliente, imputando los costos directos de transporte (combustible devengado vía GPS, peajes, tiempo de conducción y tiempo de descarga) y los costos indirectos prorrateados (bodega y administración). &  & Cap. 4.4, 7.2, 9.8, CA-11 \\
\textbf{RF-11.03} & Cálculo Diario Unificado del Indicador OTIF & Gerente Comercial (Primario) / Jefa de Calidad, Conductor repartidor propio, Conductor repartidor externo (Secundarios) & El sistema debe calcular diariamente el indicador OTIF por cliente, ruta, zona y total consolidado, definiendo «On-Time» como el cumplimiento de fecha y ventana horaria pactada e «In-Full» como la entrega del 100\% de ítems y unidades pedidas, obligando al registro de causal tipificada ante desvíos. &  & Cap. 7.1, 9.1, CA-04 \\
\textbf{RF-11.04} & Medición de Fill Rate y Perfect Order con POD Móvil & Gerente Comercial (Primario) / Conductor repartidor propio, Conductor repartidor externo (Secundarios) & El sistema debe calcular automáticamente el Fill Rate (unidades entregadas vs. pedidas) y el índice de Perfect Order contrastando la orden de venta original contra los datos de entrega física confirmados en el dispositivo móvil de reparto, sin requerir interacción tecnológica del cliente tradicional. &  & Cap. 7.1, 9.4, 16.2, CA-08 \\
\textbf{RF-11.05} & Monitoreo y Alerta de Ocupación de Flota & Planificador de rutas / Gerente Comercial & El sistema debe calcular diariamente la utilización de capacidad volumétrica (m\textasciicircum{}3) y de peso (kg) por viaje despachado, alertando aquellos camiones que salgan con una ocupación inferior al umbral configurable (ej. < 65\%). &  & Cap. 4.4, 7.2, CA-14 \\
\textbf{RF-11.06} & Tablero de Control Operacional en Tiempo Real & Gerente Comercial Gerente de Finanzas, Jefa de Calidad & El sistema debe proveer un Tablero de Control Operacional en tiempo real que consolide el avance de rutas de despacho, porcentaje de entregas cumplidas, devoluciones en tránsito, alertas activas de cadena de frío y monto recaudado en efectivo. &  & RT-05.29, Cap. 9.1, 9.6 \\
\textbf{RF-11.07} & Liquidación y Cierre Comercial por Camión & Cajero Liquidador / Conductor repartidor propio, Conductor repartidor externo, Gerente de Finanzas & El sistema debe procesar automáticamente la liquidación y cierre comercial individual por camión al retornar a base, consolidando documentos tributarios entregados, devoluciones físicas recibidas, cobranzas en efectivo y balance de envases retornables. &  & RT-05.29, CA-10 \\
\textbf{RF-11.08} & Segmentación de Clientes por Rentabilidad Neta & Gerente Comercial / Preventista, Gerente de Finanzas & El sistema debe segmentar mensualmente la cartera de clientes en matrices de rentabilidad neta (margen comercial vs. costo de servir real), sugiriendo ajustes de frecuencia de visita o umbrales mínimos de compra para clientes no rentables. &  & Cap. 7.2, 9.8, Decisión 16.1 \#7 \\
\textbf{RF-12.01} & Recepción de pedido electrónico & Cliente del canal moderno (Cadena de supermercados) & El sistema debe recibir el mensaje de pedido electrónico estructurado (EDI/EDIFACT/XML, según el estándar que exija cada cadena registrada del canal moderno). & Alta & Cap. 9.9; RT-05.23 \\
\textbf{RF-12.02} & Validación automática del pedido EDI & Cliente del canal moderno (Cadena de supermercados) & El sistema debe validar automáticamente el pedido EDI recibido, verificando estructura del mensaje, existencia del cliente, existencia de los productos, cantidades, precios y demás datos obligatorios de la cadena. & Alta & Cap. 9.9 / Cap. 17.1 \\
\textbf{RF-12.03} & Resolución del código de producto contra el maestro de Puelche & Cliente del canal moderno (Cadena de supermercados) & El sistema debe resolver el código de producto informado por la cadena contra el maestro de productos vigente de Puelche, mediante una tabla de equivalencias por cadena. & Alta & Cap. 16.1, Decisión \#11 \\
\textbf{RF-12.04} & Registro automático del pedido EDI válido en el sistema de gestión & Cliente del canal moderno (Cadena de supermercados) & El sistema debe registrar automáticamente, sin digitación manual, el pedido EDI válido (RF-12.02, RF-12.03) en el sistema de gestión empresarial. &  &  \\
\textbf{Criterio de Aceptación \#13} &  &  &  &  &  \\
\textbf{RF-12.05} & Disponibilización del pedido a la planificación de despacho & Planificador de rutas & El sistema debe poner a disposición del proceso de planificación de rutas todo pedido registrado proveniente de un canal electrónico. & Alta & Cap. 18, Criterio de Aceptación \#13 \\
\textbf{RF-12.06} & Gestión de excepciones de pedidos EDI inválidos & Operador & El sistema debe permitir gestionar, mediante una bandeja de excepciones, los pedidos que no pasaron la validación (RF-12.02) o la resolución de maestro (RF-12.03), permitiendo corregir, rechazar o reprocesar cada uno. & Alta & Cap. 17.1 / 17.2 \\
\textbf{RF-12.07} & Confirmación o rechazo del pedido a la cadena de origen & Cliente del canal moderno (Cadena de supermercados) & El sistema debe enviar a la cadena de origen una confirmación o rechazo del pedido, indicando la causa cuando no pueda ser aceptado completamente. & Alta & Cap. 9.9 \\
\textbf{RF-12.08} & Historial de estados del pedido electrónico & Operador de excepciones EDI / Planificador de rutas & El sistema debe mantener el historial de estados de cada pedido electrónico (recepción, validación, confirmación, preparación, despacho, ASN y prueba de entrega) &  &  \\
\textbf{hora y responsable/sistema,Alta,Cap. 9.1 / 9.4 / 9.9} &  &  &  &  &  \\
\textbf{RF-12.09} & Envío de aviso de despacho anticipado (ASN) & Cliente del canal moderno (Cadena de supermercados) & El sistema debe generar y enviar al cliente del canal moderno un ASN cuando su pedido queda cargado en el camión, referenciando la guía de despacho electrónica ya emitida, con productos, cantidades y hora estimada de llegada. & Alta & Cap. 9.9; entrevista Rubén Salinas \\
\textbf{RF-12.10} & Planificación de la llegada dentro de la ventana de 30 minutos & Planificador de rutas & El sistema debe planificar la llegada del camión al cliente del canal moderno dentro de la ventana horaria de 30 minutos comprometida. & Alta & Cap. 9.9 \\
\textbf{RF-12.11} & Registro del cumplimiento de la ventana de entrega & Conductor & El sistema debe registrar el cumplimiento o incumplimiento de la ventana de 30 minutos para cada entrega del canal moderno, distinguiéndola de la ventana crítica de despacho 05:30-07:00. &  &  \\
\textbf{RF-12.12} & Captura de evidencia de entrega en canal moderno & Conductor & El sistema debe capturar evidencia de entrega en el local del cliente: firma en pantalla del titular cuando pueda firmar, o fotografía de la mercadería más nombre de quien recibe cuando no sea el titular o no pueda firmar. & Alta & Cap. 9.9; Cap. 16.1, Decisión \#15 \\
\textbf{RF-12.13} & Envío de acuse de recibo digital al cliente & Cliente canal moderno & El sistema debe enviar al cliente del canal moderno un acuse de recibo digital de la entrega, articulado con la guía de despacho electrónica y su acuse de recibo tributario. & Alta & Cap. 9.9 \\
\textbf{RF-12.14} & Portal público de catálogo & Usuario no autenticado & El sistema debe publicar un portal web público, sin autenticación, con el catálogo vigente de productos (sin precios ni stock), contacto comercial e información institucional. & Media & RT-16.30 \\
\textbf{RF-12.15} & Pedido en autoservicio para clientes del canal moderno & Cliente & El portal autenticado de clientes debe permitir generar un pedido en autoservicio, validando el mismo crédito y stock que consulta el preventista. & Alta & RT-16.30 \\
\textbf{RF-12.16} & Consulta de estado de entregas en el portal de clientes & Cliente & El portal de clientes debe permitir consultar el estado de sus entregas y la ventana de entrega comprometida. & Alta & RT-16.30 / Cap. 9.1 \\
\textbf{RF-12.17} & Consulta de documentos tributarios en el portal de clientes & Cliente & El portal de clientes debe permitir consultar y descargar sus documentos tributarios (factura, guía electrónica, nota de crédito). & Alta & RT-16.30 / RT-16.14 \\
\textbf{RF-12.18} & Consulta de saldo y estado de cuenta en el portal de clientes & Cliente & El portal de clientes debe permitir consultar el saldo y estado de cuenta del cliente autenticado, con datos del sistema de gestión. & Alta & RT-16.30 \\
\textbf{RF-12.19} & Consulta de rutas asignadas en el portal de transportistas & Empresa transportista / Jefe de flota & El portal de transportistas debe permitir consultar las rutas asignadas a la empresa para el día siguiente. & Media & RT-16.30 \\
\textbf{RF-12.20} & Descarga de documentación de despacho en el portal de transportistas & Empresa transportista / Jefe de flota & El portal de transportistas debe permitir descargar la documentación de despacho de cada ruta asignada. &  &  \\
\textbf{RF-12.21} & Confirmación diaria de conductor y vehículo & Empresa transportista / Jefe de flota & El portal debe permitir confirmar, antes de la ventana de despacho, qué conductor y vehículo ejecutarán cada ruta asignada del día. & Media & Cap. 16.1, Decisión \#5 \\
\textbf{RF-12.22} & Consulta de órdenes de compra en el portal de proveedores & Proveedor & El portal de proveedores debe permitir consultar, en modo de solo lectura, el estado de sus órdenes de compra. & Media & RT-16.30 \\
\textbf{RF-12.23} & Consulta de recepciones en el portal de proveedores & Proveedor & El portal de proveedores debe permitir consultar, en modo de solo lectura, las recepciones registradas para sus órdenes. & Media & RT-16.30 \\
\textbf{RF-12.24} & Consulta de devoluciones y notas de crédito en el portal de proveedores & Proveedor & El portal de proveedores debe permitir consultar, en modo de solo lectura, sus devoluciones y notas de crédito. &  &  \\
\textbf{RF-12.25} & Registro de cliente en el portal & Cliente & El portal debe permitir al cliente del canal moderno registrarse con su RUT de empresa, correo electrónico corporativo y una contraseña, siempre que el RUT esté previamente habilitado en el sistema de gestión como cliente activo de Puelche. & Muy Alta &  \\
\textbf{RF-12.26} & Inicio de sesión en el portal de clientes & Cliente & El portal debe permitir al cliente autenticarse con su correo y contraseña registrados. & Muy Alta &  \\
\textbf{RF-12.27} & Recuperación de contraseña & Cliente & El portal debe permitir al cliente solicitar un enlace de restablecimiento de contraseña a su correo registrado cuando no recuerde su clave. & Alta &  \\
\textbf{RF-12.28} & Cierre de sesión en el portal de clientes & Cliente & El portal debe permitir al cliente cerrar su sesión de forma explícita, invalidando el token de sesión activo. & Alta &  \\
\textbf{RF-12.29} & Expiración automática de sesión por inactividad & Cliente & El portal debe cerrar automáticamente la sesión del cliente tras un período de inactividad definido, sin requerir acción explícita de cierre. & Alta &  \\
\textbf{RF-14.01} & Integración con telemetría de camiones propios & Jefe de operaciones / Gerente de operaciones & El sistema debe integrarse con la plataforma de telemetría de terceros existente que cubre los camiones propios para consumir datos de posición GPS, velocidad y kilometraje. & Muy Alta & Cap. 5, RT-17.06 \\
\textbf{RF-14.02} & Extensión de telemetrá a camiones externos & Jefe de flota / Transportista, Jefe de operaciones & El sistema debe extender la visibilidad de telemetría a los camiones de transportistas externos contratados cuyo operador cuente con GPS habilitado, mediante un mecanismo acordado mediante la aplicacion. & Alta & Cap. 2.3, Cap. 17.4 (punto 9). \\
\textbf{RF-14.03} & Visualización de rutas planificadas vs reales & Gerente de operaciones, Planificador de rutas & El sistema debe mostrar en el mapa las rutas planificadas vs las rutas reales recorridas por cada camión, para identificar desviaciones y optimizar las rutas futuras. & Alta & Cap. 3 y entrevista de nelson \\
\textbf{RF-14.04} & Alerta de desviación de ruta & Centro de control / Jefe de operaciones & El sistema debe generar una alerta cuando el camión se desbie más de 3 km de la ruta planificada, para que el centro de control pueda reaccionar. & Alta & Cap. 9.1 \\
\textbf{RF-14.05} &  & Conductor repartidor, Jefe de flota & El sistema debe registrar el tiempo de conducción y descanso del conductor, utilizando los datos de posicionamiento (con la debida consideración de la objeción sindical) para cumplir con el régimen especial de jornada. & Media & Cap. 12 (Régimen jornada conductores), Cap. 16.1 (Decisión \#13), Restricción \#10 \\
\textbf{RF-14.06} & Integración de telemetría con costo de servir & Gerente de finanzas, Jefe de operaciones & El sistema debe integrar los datos de telemetría con el módulo de costo de servir, imputando kilómetros recorridos y tiempo de viaje a cada entrega y cliente. & Muy Alta & Cap. 9.8 \\
\textbf{RF-14.07} & Vinculación conductor-camión en cada viaje & Conductor repartidor, Jefe de flota & El sistema debe registrar qué conductor propio o de transportista externo está asociado a cada camión en cada viaje, consumiendo la identidad validada por el mecanismo de autenticación (RF-15.02). Esta vinculación es necesaria para imputar responsabilidades en cobros, entregas y eventos de ruta, y para resolver la Decisión \#5. & Alta & Cap. 16.1 (Decisión \#5), Cap. 4.6 \\
\textbf{RF-14.08} & Gestión de geocercas para control de llegada/salida & Jefe de operaciones / Gerente de operaciones, Conductor repartidor & El sistema debe permitir la gestión de geocercas zonas geográficas, para identificar cuándo un camión entra o sale de una zona ej: CD Talca, zona de clientes, para medir tiempos de llegada y salida. & Media & Cap. 9.9 y entrevista con ruben salinas \\
\textbf{RF-14.09} & Reportes de kilometraje y consumo de combustible & Gerente de finanzas, Jefe de operaciones & El sistema debe consolidar los datos de telemetría de todos camiones propios para generar reportes de kilometraje total y consumo de combustible estimado, apoyando la gestión de costos de transporte. & Alta & Cap. 7.2 (costo desconocido). \\
\end{longtable}

